\section{Related Work}

\paragraph{Metacell methods.}
Metacell approaches address the sparsity and noise in single-cell data by aggregating highly similar cells into stable representative profiles.
The standard pipeline constructs a cell--cell affinity graph---typically a $k$-nearest-neighbor graph with kernel-based similarity on PCA embeddings---and then partitions it into coherent groups.
MetaCell~\citep{metacell} introduced this paradigm by partitioning a kNN graph into balanced, densely connected groups using community detection algorithms~\citep{louvain, traag2019leiden} that operate on graph topology alone, without incorporating gene expression features directly.
SEACells~\citep{seacells} refines the affinity with an adaptive Gaussian kernel that accounts for local density variations, and applies archetypal analysis~\citep{cutler1994archetypal} to identify extreme points on the data manifold as representative cell groups.
Other methods such as Metacell-2~\citep{metacell2} and SuperCell~\citep{supercell} similarly aggregate cells based on transcriptomic graph structure, but like MetaCell, they remain restricted to within-batch similarity.

These methods share two limitations.
First, because batch effects induce systematic shifts across datasets, the affinity graph is reliable only within individual batches, and metacells remain batch-specific even when the underlying cell states are shared across datasets.
Second, groupings are defined purely by affinity structure, with no learning objective that validates whether the grouped cells share coherent gene expression patterns.
MetaQ~\citep{metaq} attempts to address the second limitation by replacing graph construction with end-to-end representation learning via a vector-quantization objective~\citep{vqvae}, but without an explicit batch correction mechanism the model can reduce its loss by learning batch-specific metacells---an always-available signal---rather than shared biological states.
SURE~\citep{sure} takes a generative approach that jointly constructs metacells and removes batch effects, but like MetaQ, optimizes reconstruction without affinity guidance, failing to preserve local cell--cell similarity structure---a limitation shared more broadly by dedicated batch effect correction methods.

\paragraph{Batch effect correction.}
Other VAE-based methods such as scVI~\citep{scvi} similarly use conditional variational autoencoders~\citep{cvae} to disentangle biological variation from batch effects, but optimize a global reconstruction loss without affinity guidance, causing rare cell states to be absorbed into larger clusters.
scPoli~\citep{scpoli} partially addresses this by anchoring prototypes to known cell types, though this requires accurate annotations, which are often incomplete or inconsistent in practice.
Non-VAE approaches such as Harmony~\citep{harmony} and BBKNN~\citep{bbknn} take different strategies---iterative embedding alignment and graph-level connectivity enforcement, respectively---but introduce their own distortions: Harmony can disrupt within-batch cell-type structure, while BBKNN assumes all cell types are present in every batch.
Taken together, none of these methods incorporate an affinity graph during training---they align cells but do not explicitly preserve the community structure defined by cell--cell relationships within each dataset.
Preserving such structure requires objectives that operate at the level of groups rather than individual cell distances, which connects to the literature on deep clustering.

\paragraph{Deep clustering and prototype-based learning.}
Deep clustering methods~\citep{deepcluster} jointly learn embeddings and cluster assignments, alternating between assigning cells to clusters and updating the network to predict those assignments.
SeLa~\citep{asano2019self} replaces k-means with Sinkhorn--Knopp~\citep{sinkhorn} assignment for balanced pseudo-labels; SwAV~\citep{swav} extends this online with learnable prototypes.
However, this enforces uniform usage across prototypes---which does not reflect the highly variable cluster sizes in biological data---and defines positive structure through data augmentation rather than input-defined similarity, providing no information about how similar two cells actually are; in single-cell settings, a cell--cell affinity graph provides a more principled source of structure.
Graph-based methods such as Parametric UMAP~\citep{parametric_umap} do incorporate such a graph, learning embeddings that preserve pairwise affinities, but target distances between individual cells rather than cluster-level assignments---in metacell learning, what matters is which cells are grouped together, not the precise geometry of the embedding space.
A prior version of this work~\citep{golpayegani2025interpretable} adapted SwAV-style prototype learning to single-cell transcriptomics using an affinity graph, demonstrating that shared prototypes can yield interpretable representations, but shares the same limitation of uniform prototype usage.
A related failure mode across prototype-based methods is prototype collapse, where a subset of prototypes captures most cells and the rest go unused. Arteaga et al.~\citep{arteaga2026prototypes} address this via decoupled prototype updates modelled as a Gaussian mixture, though without enforcing distinctness, prototypes may remain redundant rather than capturing distinct clusters.

Existing approaches thus leave a gap: affinity-preserving methods remain restricted to within-batch aggregation and rely solely on graph topology without incorporating feature information, preventing a balance between affinity-based and expression-based grouping; methods that align cells across batches do not preserve affinity structure; and prototype-based methods that could bridge both enforce uniform prototype usage via balanced assignment, limiting their ability to capture biological communities of varying sizes.
